\documentclass[11pt,a4paper]{article}

\usepackage[utf8]{inputenc}
\usepackage[T1]{fontenc}
\usepackage[french]{babel}
\usepackage[margin=2.3cm]{geometry}
\usepackage{amsmath,amssymb}
\usepackage{booktabs}
\usepackage{tabularx}
\usepackage{enumitem}
\usepackage[hidelinks]{hyperref}

\newcolumntype{C}{>{\centering\arraybackslash}X}
\setlist{itemsep=2pt,topsep=4pt}

\title{Projet SDP -- Jalon 3\\[2pt]\large Attaquer l'apprentissage d'un modèle MR-Sort par injection de données}
\author{Mohamed Salh \and Mohamed Chebli \and Quentin Robidou}
\date{Novembre 2026}

\begin{document}
\maketitle

\setcounter{section}{-1}
% =====================================================================
\section{Cadrage du problème}

\subsection{Les acteurs}

Un utilisateur A envoie à un utilisateur B un ensemble de données $(x, k)$ classées par son modèle MR-Sort,
afin que B apprenne ce modèle ou un modèle équivalent : c'est un problème inv-MR-Sort. Un attaquant C cherche
à perturber cet apprentissage en ajoutant ses propres données, synthétiques, à celles de A. Nous nous plaçons
du point de vue de C : il s'agit de concevoir la meilleure attaque possible.

\subsection{Les hypothèses à faire varier}

Le sujet nous laisse cadrer le problème. Dans tous les scénarios, B reçoit en une fois l'ensemble des données de
A et de C, mélangées et sans date d'arrivée. Nous retenons trois hypothèses, qui définissent les scénarios :
\begin{enumerate}
  \item \textbf{B se défend-il ?} Soit B apprend naïvement sur toutes les données reçues, soit il cherche à
        repérer et écarter les données de C.
  \item \textbf{C a-t-il accès aux données envoyées par A ?} Uniquement aux données, jamais au modèle de A.
  \item \textbf{Quelle part des données reçues par B provient de C ?} Au plus 50\,\%.
\end{enumerate}

\subsection{Ce que sait l'attaquant}

C ne connaît que le contour du problème : le nombre de critères $n$, le nombre de catégories $p$, le fait que
les évaluations sont dans $[0,1]^n$ et que les critères sont à maximiser. Il ne connaît ni la façon dont le
modèle de A a été choisi, ni la procédure d'apprentissage de B, ni la façon dont B se défend. Pour construire
une attaque forte contre un MR-Sort quelconque, C a cependant besoin, dans certaines stratégies, de simuler un
apprentissage : il simule alors une procédure inv-MR-Sort générique, offrant un minimum de robustesse, et nous
prenons lincs (section~1). Ce n'est pas une hypothèse sur B, mais un manque d'alternative : C ne peut pas
simuler la vraie stratégie de B, qu'il ne connaît pas, mais il peut simuler une procédure inv-MR-Sort générique.

\subsection{Mesure de la qualité d'une attaque}

Une attaque est d'autant meilleure que le modèle appris par B est éloigné de celui de A. Comme vu au jalon 1,
deux modèles aux paramètres différents peuvent être équivalents : nous mesurons donc leur proximité de façon
empirique, par la proportion de points sur lesquels ils prennent la même décision. Le score d'une stratégie
s'obtient ainsi :
\begin{enumerate}
  \item tirer un modèle MR-Sort A au hasard, et $m_{\text{train}}$ objets uniformes dans $[0,1]^n$ classés par A ;
  \item C ajoute ses données ; si une part $\alpha$ des données reçues par B doit provenir de C, C ajoute
        $m_{\text{train}} \cdot \alpha / (1-\alpha)$ objets (B s'entraîne donc sur plus de $m_{\text{train}}$ objets) ;
  \item B apprend un modèle sur l'ensemble, remélangé ;
  \item sur $m_{\text{test}}$ objets uniformes nouveaux, mesurer la proportion d'objets affectés à la même
        catégorie par A et par B (l'\emph{accord}) ;
  \item répéter sur plusieurs modèles A et faire la moyenne : le score de la stratégie est cet accord moyen,
        d'autant meilleur pour C qu'il est bas.
\end{enumerate}
Chaque couple (stratégie, proportion) est évalué sur \textbf{ses propres modèles A}, tirés indépendamment
(60 par couple, avec $n$ entre 3 et 6 et $p$ entre 2 et 4 tirés à chaque fois et communiqués à C,
$m_{\text{train}} = 300$, $m_{\text{test}} = 10\,000$). Chaque score est ainsi une estimation indépendante de
l'efficacité moyenne de la stratégie à cette proportion, que nous accompagnons de son intervalle de confiance à
95\,\%, et que nous comparons à une \textbf{référence sans attaque} ($\alpha = 0$), elle aussi estimée sur 60
modèles A.

\subsection{Évaluer exactement le modèle appris}

Le modèle appris par B doit être évalué tel que l'algorithme d'apprentissage le définit. lincs calcule en
simple précision (nombres à virgule sur 32 bits, environ 7 chiffres significatifs) : il compare les évaluations
aux profils et additionne les poids critère par critère dans ce format, et déclare une coalition suffisante si
la somme atteint 1. Recalculer la concordance en double précision (64 bits) ne donne pas toujours le même
modèle : trois poids de $0{,}33333299$ font $0{,}99999997 < 1$ en double précision, alors qu'en simple
précision cette somme est arrondie à 1, le nombre représentable le plus proche. Sur certains modèles, la
coalition change donc de statut, et le modèle évalué n'est plus celui que lincs a appris (dans un cas sans
attaque, l'accord tombait ainsi à 62\,\% au lieu de 97\,\%). Nous reproduisons donc exactement le calcul de
lincs ; sur 400\,000 points de test, notre évaluation et celle de lincs ne divergent sur aucun.

% =====================================================================
\section{Premier scénario : attaque à l'aveugle d'un apprentissage naïf}

Ce premier scénario est le plus simple : C attaque \textbf{à l'aveugle}
(sans accès aux données de A), et B \textbf{ne se défend pas}. Nous testons les parts de données de C
suivantes : 10\,\%, 25\,\%, $1/3$ et 50\,\%.

\subsection{Choix de l'algorithme d'apprentissage de B}

Il faut d'abord fixer ce que signifie « B ne se défend pas ». Nous nous plaçons du point de vue d'un B naïf,
qui ne s'attend ni à être attaqué ni à recevoir des données bruitées. Deux algorithmes viennent naturellement
de nos travaux précédents, et B pourrait préférer l'un ou l'autre pour des raisons sans rapport avec
l'attaque :
\begin{itemize}
  \item le \textbf{programme du jalon 1} (question vii), résolu avec Gurobi, qui maximise la marge $\alpha$
        entre la concordance de chaque objet et le seuil $\lambda$ (limité ici à 30\,s) ;
  \item l'\textbf{heuristique de lincs} utilisée au jalon 2, qui cherche le modèle classant correctement le plus
        grand nombre d'objets (arrêt quand tout est bien classé, après 20 itérations sans progrès, ou après 30\,s).
\end{itemize}
Il reste à choisir lequel rend l'expérience la plus intéressante. A priori, lincs semblait le moins
intéressant : maximiser le nombre d'objets bien classés le rend plus robuste, ce qui risquait d'atténuer
l'effet des attaques, et en particulier de désavantager injustement les attaques par données aléatoires, que
lincs peut simplement sacrifier.

L'expérience a montré l'inverse. Le programme Gurobi est si peu robuste qu'il s'effondre à la moindre attaque.
Dès que les données de A et de C sont incompatibles, aucune marge positive n'est possible, et la marge optimale
vaut $\alpha^* = 0$ : elle est atteinte par une solution triviale où toutes les concordances sont égales au
seuil, ce qui place tous les objets dans la meilleure catégorie. Gurobi infère donc un modèle nul, quelle que
soit l'attaque :

\begin{table}[h]
\centering
\begin{tabularx}{\textwidth}{>{\raggedright\arraybackslash}p{5.2cm}CCCCC}
\toprule
Part des données de C & 0\,\% & 10\,\% & 25\,\% & $1/3$ & 50\,\% \\
\midrule
lincs, données aléatoires   & 99,2\,\% & 98,3\,\% & 98,2\,\% & 97,9\,\% & 98,9\,\% \\
Gurobi, données aléatoires  & 99,3\,\% & 8,2\,\%  & 8,2\,\%  & 8,2\,\%  & 8,2\,\% \\
Gurobi, autre modèle MR-Sort & 99,1\,\% & 8,2\,\% & 8,2\,\%  & 8,2\,\%  & 8,2\,\% \\
\bottomrule
\end{tabularx}
\caption{Accord entre B et A pour un même modèle A ($n=6$, $p=3$). Avec Gurobi, toutes les attaques donnent
exactement le même modèle dégénéré.}
\label{tab:algo}
\end{table}

Avec Gurobi, toutes les stratégies produisent le même résultat, l'un des pires possibles : il devient
impossible de déterminer quelle stratégie d'attaque est la meilleure. Imposer $\lambda \ge 0{,}5$, comme dans
l'article de Leroy, Mousseau et Pirlot, ne suffit pas (accord moyen de 14 à 20\,\% sur 30 modèles), car
d'autres solutions triviales de marge nulle existent. Nous retenons donc \textbf{lincs}, seul scénario à la
fois crédible et intéressant.

\subsection{Stratégies d'attaque}

Sans connaître A, C peut seulement choisir où placer ses points et comment les étiqueter. Nous comparons six
stratégies :
\begin{itemize}
  \item \textbf{S1 -- Bruit aléatoire} : points uniformes, catégories tirées au hasard.
  \item \textbf{S2 -- Autre modèle MR-Sort} : points uniformes, classés par un MR-Sort que C tire au hasard.
  \item \textbf{S3 -- Catégorie extrême} : points uniformes, tous dans la meilleure catégorie, ou tous dans la
        pire (choix tiré au hasard). Ces données sont représentables par un MR-Sort, mais contredisent A
        presque partout.
  \item \textbf{S4 -- Dictature d'un critère} : points uniformes, classés par un MR-Sort où un seul critère, tiré
        au hasard, décide de tout. Elle vise les poids plutôt que les profils.
  \item \textbf{S5 -- Anti-monotone} : points uniformes, classés à l'envers d'un MR-Sort aléatoire (catégorie
        $p-1-k$). Aucun MR-Sort ne peut expliquer ces données.
  \item \textbf{S6 -- Catégorie extrême au centre} : comme S3, mais avec des points dans $[0{,}25\,;0{,}75]^n$.
        Près des coins de l'espace, la catégorie de n'importe quel MR-Sort est presque toujours la pire ou la
        meilleure : seuls les points intérieurs ont une catégorie incertaine, et peuvent déplacer les profils.
\end{itemize}

\subsection{Résultats}

Sans attaque, B obtient $97{,}5 \pm 1{,}0$\,\% d'accord en moyenne avec A (98,5\,\% en médiane ; lincs converge
mal sur 2 modèles sur 60, à 76\,\%).

\begin{table}[h]
\centering
\begin{tabularx}{\textwidth}{>{\raggedright\arraybackslash}p{4.6cm}CCCC}
\toprule
Stratégie & 10\,\% & 25\,\% & $1/3$ & 50\,\% \\
\midrule
S1 -- Bruit aléatoire         & $98{,}1 \pm 0{,}4$ & $96{,}6 \pm 0{,}8$ & $95{,}5 \pm 1{,}0$ & $94{,}5 \pm 1{,}3$ \\
S2 -- Autre modèle MR-Sort    & $97{,}5 \pm 0{,}5$ & $96{,}2 \pm 0{,}8$ & $94{,}0 \pm 1{,}3$ & $72{,}2 \pm 4{,}5$ \\
S3 -- Catégorie extrême       & $97{,}3 \pm 0{,}7$ & $95{,}3 \pm 1{,}0$ & $92{,}3 \pm 1{,}7$ & $\mathbf{65{,}4 \pm 6{,}0}$ \\
S4 -- Dictature d'un critère  & $97{,}4 \pm 0{,}6$ & $95{,}7 \pm 1{,}0$ & $92{,}5 \pm 1{,}4$ & $\mathbf{68{,}1 \pm 4{,}0}$ \\
S5 -- Anti-monotone           & $97{,}7 \pm 0{,}5$ & $96{,}4 \pm 0{,}8$ & $94{,}5 \pm 1{,}2$ & $85{,}9 \pm 3{,}8$ \\
S6 -- Extrême au centre       & $97{,}3 \pm 0{,}7$ & $\mathbf{88{,}3 \pm 3{,}0}$ & $\mathbf{79{,}7 \pm 3{,}7}$ & $70{,}6 \pm 4{,}2$ \\
\bottomrule
\end{tabularx}
\caption{Accord moyen (\%) entre B et A selon la part des données de C, avec l'intervalle de confiance à 95\,\%
(60 modèles A tirés indépendamment pour chaque case ; sans attaque : $97{,}5 \pm 1{,}0$\,\%).}
\label{tab:resultats}
\end{table}

Le tableau~\ref{tab:resultats} conduit aux constats suivants :
\begin{itemize}
  \item \textbf{À 10\,\%, aucune attaque n'est efficace} : tous les scores restent dans l'intervalle de confiance de
        la référence sans attaque.
  \item \textbf{À 25\,\% et $1/3$, S6 est la meilleure stratégie}, et c'est la seule dont l'intervalle de confiance se
        détache nettement des autres ($88{,}3 \pm 3{,}0$ puis $79{,}7 \pm 3{,}7$\,\%). En concentrant ses points au centre
        de l'espace, C devient majoritaire dans cette zone, alors qu'il reste minoritaire sur l'ensemble des
        données : lincs, qui maximise le nombre d'objets bien classés, préfère céder sur cette zone. C'est la
        \textbf{proportion locale} des données de C qui compte, plus que leur proportion globale.
  \item \textbf{À 50\,\%, S3 et S4 sont les meilleures stratégies} ($65{,}4 \pm 6{,}0$ et $68{,}1 \pm 4{,}0$\,\%), sans
        qu'on puisse les départager, ni les séparer nettement de S6 et S2 ($70{,}6 \pm 4{,}2$ et $72{,}2 \pm 4{,}5$\,\%) :
        les intervalles se recouvrent. À parts égales, C n'a plus besoin de se concentrer : des données cohérentes
        et très éloignées de A sur tout l'espace suffisent.
  \item \textbf{Une attaque efficace doit être représentable par un MR-Sort.} Le bruit (S1) n'a presque aucun effet,
        même à 50\,\% ($94{,}5 \pm 1{,}3$\,\%), et les données anti-monotones (S5) sont nettement moins
        efficaces qu'un MR-Sort aléatoire (S2) à 50\,\% : lincs écarte facilement des données qu'aucun MR-Sort ne peut
        expliquer. Les stratégies les plus efficaces (S3, S4, S6) envoient au contraire des données cohérentes
        mais très éloignées d'un modèle typique.
  \item \textbf{Les problèmes à deux catégories résistent mieux} : à 50\,\%, l'accord vaut 70\,\% pour $p=2$ contre 61
        à 65\,\% pour $p=3$ ou 4 avec S3, et 80\,\% contre 56 à 65\,\% avec S4.
\end{itemize}
Deux réserves : les intervalles restent larges à 50\,\%, car l'effet de l'attaque varie beaucoup d'un modèle A à
l'autre ; et S3 et S6 choisissent au hasard la meilleure ou la pire catégorie, sans que nous sachions encore si
l'un des deux choix est plus efficace.

% =====================================================================
\section{Deuxième scénario : l'attaquant a accès aux données de A}

Le cadre reste le même (B naïf qui apprend avec lincs), mais C reçoit désormais
toutes les données envoyées par A, sans connaître son modèle. Il peut en inférer un modèle $\hat{A}$ : les
données de A n'étant pas bruitées, $\hat{A}$ appris avec lincs est très proche de A (98,3\,\% d'accord en
moyenne). C choisit aussi un sens de décalage des catégories ($+1$ ou $-1$), celui qui laisse le moins
d'étiquettes de $\hat{A}$ bloquées à une extrémité.

Pour l'une des stratégies, C doit simuler un apprentissage. Il ne connaît pas la procédure de B : il simule une
procédure inv-MR-Sort générique, offrant un minimum de robustesse, et nous prenons lincs. L'objectif reste de
construire un ensemble de données trompeur pour toute procédure de type inv-MR-Sort. Dans ce scénario, B ne se
défend pas, et apprend lui aussi avec une procédure inv-MR-Sort générique : lincs, retenu en partie~1 pour des
raisons similaires. Ce choix commun n'est pas fait pour avantager C, il vient du manque d'alternative ; mais il
faut garder à l'esprit qu'à l'évaluation, l'attaque en est avantagée, puisque C simule exactement la procédure de
B. Ce n'est plus le cas dans la partie~4, où B se défend.

\subsection{Stratégies d'attaque}

\begin{itemize}
  \item \textbf{S7 -- Modèle $\hat{A}$ décalé} : points uniformes, étiquetés par $\hat{A}$ décalé d'une catégorie. Ce
        modèle décalé est encore un MR-Sort, donc représentable, et il contredit $\hat{A}$ sur tous les objets qui
        ne sont pas déjà à une extrémité.
  \item \textbf{S8 -- Copie à étiquette décalée} : C reprend les points mêmes de A et décale leur étiquette d'une
        catégorie (dans l'autre sens si elle est déjà à une extrémité). Le conflit est maximal : ce sont les
        mêmes objets, avec deux étiquettes différentes.
  \item \textbf{S9 -- Viser les frontières} : C tire ses points autour des profils de $\hat{A}$, là où la catégorie
        bascule, et les étiquette du côté opposé à celui que $\hat{A}$ leur attribue.
  \item \textbf{S10 -- Optimisation du placement} : l'attaque idéale consisterait à optimiser directement
        chaque point de C, ses $n$ coordonnées et son étiquette, pour minimiser l'accord entre A et le modèle
        appris par B. Avec $\gamma\,m_{\text{train}}(n+1)$ paramètres (jusqu'à 2\,100) et un objectif dont chaque
        évaluation demande un apprentissage, ce problème est hors de portée. Nous le simplifions en optimisant la
        \emph{distribution} des points : C les place uniformément dans une boîte $[c-r,\,c+r]$ et les étiquette par
        $\hat{A}$ décalé ; les paramètres sont le centre $c$, la demi-largeur $r$ et le sens du décalage, soit
        $n+2$ paramètres. Un algorithme génétique (10 individus, 7 générations) les optimise ; chaque individu est
        évalué en simulant l'apprentissage avec lincs et en mesurant l'accord entre le modèle obtenu et $\hat{A}$, avec de
        nouveaux points, une nouvelle graine et un nouvel échantillon d'évaluation à chaque fois, pour ne pas
        sur-optimiser contre une exécution particulière. L'accord n'étant pas une fonction dérivable des
        paramètres (fonction en escalier, apprentissage en boîte noire), une descente de gradient n'est pas
        possible. Cette stratégie, plus coûteuse (environ 1 min 20 par modèle A en médiane), n'a été
        évaluée que sur 10 modèles A par proportion, au lieu de 60.
\end{itemize}

\subsection{Résultats}

\begin{table}[h]
\centering
\begin{tabularx}{\textwidth}{>{\raggedright\arraybackslash}p{5.0cm}CCCC}
\toprule
Stratégie & 10\,\% & 25\,\% & $1/3$ & 50\,\% \\
\midrule
S3 -- Catégorie extrême (aveugle)   & $97{,}3 \pm 0{,}7$ & $95{,}3 \pm 1{,}0$ & $92{,}3 \pm 1{,}7$ & $65{,}4 \pm 6{,}0$ \\
S6 -- Extrême au centre (aveugle)   & $97{,}3 \pm 0{,}7$ & $88{,}3 \pm 3{,}0$ & $79{,}7 \pm 3{,}7$ & $70{,}6 \pm 4{,}2$ \\
\midrule
S7 -- Modèle $\hat{A}$ décalé       & $96{,}9 \pm 0{,}6$ & $91{,}9 \pm 1{,}4$ & $89{,}8 \pm 1{,}9$ & $55{,}0 \pm 2{,}7$ \\
S8 -- Copie à étiquette décalée     & $97{,}6 \pm 0{,}5$ & $94{,}8 \pm 1{,}0$ & $94{,}2 \pm 1{,}0$ & $51{,}5 \pm 4{,}1$ \\
S9 -- Viser les frontières          & $96{,}2 \pm 0{,}6$ & $91{,}7 \pm 1{,}2$ & $87{,}5 \pm 1{,}6$ & $79{,}3 \pm 2{,}3$ \\
S10 -- Optimisation (10 modèles)    & $\mathbf{88{,}5 \pm 3{,}1}$ & $\mathbf{69{,}9 \pm 8{,}1}$ & $\mathbf{62{,}1 \pm 9{,}1}$ & $\mathbf{23{,}1 \pm 13{,}8}$ \\
\bottomrule
\end{tabularx}
\caption{Accord moyen (\%) entre B et A selon la part des données de C, avec l'intervalle de confiance à 95\,\%
(60 modèles A tirés indépendamment par case, 10 pour S10 ; sans attaque : $97{,}5 \pm 1{,}0$\,\%). Les stratégies à
l'aveugle S3 et S6 sont rappelées de la partie~1.}
\label{tab:informed}
\end{table}

Le tableau~\ref{tab:informed} conduit aux constats suivants :
\begin{itemize}
  \item \textbf{À 50\,\%, connaître les données de A donne un avantage net.} La copie à étiquette décalée et le
        modèle $\hat{A}$ décalé font tomber l'accord à $51{,}5 \pm 4{,}1$ et $55{,}0 \pm 2{,}7$\,\%, nettement en
        dessous de la meilleure stratégie à l'aveugle (catégorie extrême, $65{,}4 \pm 6{,}0$\,\%) : contredire
        exactement A sur tout l'espace est plus efficace que de contredire un modèle typique.
  \item \textbf{En dessous de 50\,\%, les stratégies simples ne suffisent pas.} À 25\,\% et $1/3$, S7, S8 et S9 font
        moins bien que S6, qui n'utilise pas les données de A : C reste minoritaire partout, et les contradictions
        qu'il répartit sur tout l'espace sont écartées par lincs. Seule la concentration des points (S6) paie.
  \item \textbf{Ce qui paie le plus, c'est de pouvoir simuler l'apprentissage.} L'optimisation S10 est la meilleure
        stratégie à toutes les proportions, avec des intervalles séparés de ceux des autres stratégies, même sur
        10 modèles. C'est la seule à agir dès 10\,\% ($88{,}5 \pm 3{,}1$\,\%, contre 96 à 98\,\% pour les autres), et à
        50\,\% elle fait tomber l'accord à $23{,}1 \pm 13{,}8$\,\% : B apprend un modèle en désaccord avec A sur plus
        des trois quarts de l'espace en moyenne. Elle ne cherche pourtant qu'une dizaine de paramètres, avec 80
        évaluations.
  \item \textbf{La proportion locale reste la clé.} S10 ne fait que choisir où concentrer les points de C et dans
        quel sens décaler les étiquettes : c'est la généralisation de l'idée de S6, placée là où elle fait le plus
        de dégâts pour le modèle A attaqué.
\end{itemize}
Les résultats de S10 ne portent que sur 10 modèles par proportion et un seul lancement de l'algorithme génétique
par modèle : les intervalles de confiance sont larges. Sur un modèle par proportion, l'optimisation échoue
presque complètement (accord de 98\,\% à 25\,\% et à $1/3$, de 75\,\% à 50\,\%), ce qui élargit ces intervalles.

% =====================================================================
\section{Troisième scénario : B se défend}

B sait désormais qu'un attaquant peut injecter des données, sans savoir s'il y en a un, ni quelle stratégie il
suit, ni dans quelle proportion. Nous nous plaçons dans le cadre de la partie~1 : C attaque à l'aveugle avec
l'une des stratégies S1 à S6. Nous procédons en trois phases : des défenses spécialisées, chacune conçue contre
une famille d'attaques (phase~1) ; une défense hybride, qui détecte l'attaque puis applique la défense adaptée
(phase~2) ; enfin la meilleure réponse de C face à ces défenses (phase~3).

\subsection{Défenses spécialisées}

\begin{itemize}
  \item \textbf{D0 -- Aucune défense} (contre S1) : B apprend avec lincs, comme dans les parties 1 et 2. lincs
        maximise le nombre d'objets bien classés, ce qui suffit à écarter le bruit.
  \item \textbf{D1 -- Plus grand segment} (contre S2) : si C envoie des données classées par un autre MR-Sort, le
        jeu reçu est un mélange de deux MR-Sort. B apprend deux modèles avec l'heuristique du jalon~2 (départ
        « modèle unique », 30\,s par apprentissage) et garde celui qui explique le plus d'objets, en supposant que
        A est majoritaire.
  \item \textbf{D2 -- Nettoyage de la monotonie} (contre S3, S5 et S6) : un objet qui en domine un autre sur
        tous les critères ne peut pas avoir une catégorie inférieure dans un MR-Sort. B retire, un par un, l'objet
        impliqué dans le plus de violations de cette règle, jusqu'à ce qu'il n'en reste plus, puis apprend avec
        lincs sur les objets restants.
  \item \textbf{D3 -- Rejet du dictateur} (contre S4) : B apprend deux modèles comme pour D1 ; si un seul des deux
        est dictatorial (un critère suffit à lui seul à atteindre le seuil de majorité), il garde l'autre, sinon il
        garde celui du plus grand segment.
\end{itemize}
Les modèles appris par D1 et D3 sont évalués exactement comme ceux de lincs (section~0), y compris pendant
l'heuristique à deux modèles, quand elle réaffecte les objets aux segments.

\subsection{Phase 1 : chaque défense contre chaque attaque}

Chaque défense est évaluée sans attaque, et contre chaque stratégie à $1/3$ et 50\,\%, les deux proportions où
les attaques de la partie~1 ont un effet. Faute de temps (l'heuristique à deux modèles prend environ 30\,s par
jeu), chaque case n'a que 20 modèles A, tirés indépendamment.

\begin{table}[h]
\centering
\small
\begin{tabularx}{\textwidth}{>{\raggedright\arraybackslash}p{3.3cm}cCCCC}
\toprule
Attaque & C & D0 -- aucune & D2 -- monotonie & D1 -- plus grand segment & D3 -- rejet du dictateur \\
\midrule
Sans attaque         & 0       & $\mathbf{98{,}6 \pm 0{,}3}$ & $\mathbf{98{,}6 \pm 0{,}4}$ & $81{,}4 \pm 5{,}6$ & $83{,}5 \pm 8{,}0$ \\
\midrule
S1 -- Bruit          & $1/3$   & $94{,}6 \pm 1{,}8$ & $\mathbf{96{,}7 \pm 1{,}1}$ & $92{,}6 \pm 4{,}7$ & $86{,}1 \pm 10{,}4$ \\
                     & 50\,\%  & $\mathbf{95{,}1 \pm 2{,}5}$ & $93{,}3 \pm 2{,}0$ & $91{,}5 \pm 3{,}8$ & $86{,}1 \pm 8{,}0$ \\
S2 -- Autre MR-Sort  & $1/3$   & $\mathbf{95{,}4 \pm 2{,}2}$ & $94{,}7 \pm 1{,}6$ & $86{,}0 \pm 5{,}9$ & $89{,}6 \pm 8{,}1$ \\
                     & 50\,\%  & $\mathbf{81{,}8 \pm 4{,}3}$ & $74{,}0 \pm 4{,}6$ & $62{,}6 \pm 9{,}3$ & $70{,}4 \pm 9{,}9$ \\
S3 -- Extrême        & $1/3$   & $\mathbf{95{,}2 \pm 2{,}8}$ & $90{,}9 \pm 3{,}3$ & $74{,}3 \pm 11{,}1$ & $85{,}5 \pm 8{,}1$ \\
                     & 50\,\%  & $61{,}3 \pm 11{,}5$ & $\mathbf{70{,}0 \pm 7{,}8}$ & $52{,}7 \pm 16{,}2$ & $66{,}9 \pm 10{,}2$ \\
S4 -- Dictature      & $1/3$   & $\mathbf{92{,}0 \pm 2{,}7}$ & $91{,}5 \pm 3{,}1$ & $85{,}7 \pm 8{,}8$ & $86{,}0 \pm 7{,}0$ \\
                     & 50\,\%  & $67{,}3 \pm 7{,}7$ & $65{,}1 \pm 5{,}5$ & $67{,}1 \pm 10{,}6$ & $\mathbf{70{,}0 \pm 6{,}5}$ \\
S5 -- Anti-monotone  & $1/3$   & $93{,}0 \pm 2{,}5$ & $\mathbf{94{,}2 \pm 2{,}2}$ & $92{,}9 \pm 4{,}8$ & $91{,}2 \pm 4{,}5$ \\
                     & 50\,\%  & $82{,}8 \pm 5{,}5$ & $76{,}3 \pm 8{,}2$ & $83{,}7 \pm 11{,}8$ & $\mathbf{87{,}8 \pm 7{,}7}$ \\
S6 -- Extrême centre & $1/3$   & $78{,}0 \pm 4{,}9$ & $72{,}4 \pm 8{,}8$ & $\mathbf{83{,}1 \pm 8{,}1}$ & $81{,}6 \pm 9{,}8$ \\
                     & 50\,\%  & $70{,}5 \pm 8{,}2$ & $\mathbf{77{,}0 \pm 7{,}3}$ & $51{,}6 \pm 11{,}4$ & $64{,}4 \pm 11{,}9$ \\
\midrule
Pire cas sous attaque &        & $61{,}3$ & $\mathbf{65{,}1}$ & $51{,}6$ & $64{,}4$ \\
\bottomrule
\end{tabularx}
\caption{Accord moyen (\%) entre B et A pour chaque défense, avec l'intervalle de confiance à 95\,\% (20 modèles A
tirés indépendamment pour chaque case). En gras, la meilleure défense de chaque ligne.}
\label{tab:defenses}
\end{table}

Le tableau~\ref{tab:defenses} conduit aux constats suivants :
\begin{itemize}
  \item \textbf{Le nettoyage de la monotonie (D2) est la meilleure défense isolée.} Il ne coûte rien sans attaque
        ($98{,}6$\,\%, comme D0), et c'est la défense dont le pire cas est le meilleur ($65{,}1$\,\%, contre S4 à
        50\,\%). Son avance sur D3 ($64{,}4$\,\%) n'est toutefois pas significative.
  \item \textbf{Ne rien faire reste souvent le meilleur choix.} lincs seul est le meilleur, ou à égalité, dans 5 cas
        sur 12, en particulier à $1/3$ : à cette proportion, seule S6 le fait descendre sous 90\,\%.
  \item \textbf{Les défenses à deux modèles coûtent cher sans attaque} : environ 15 points ($81{,}4$ et
        $83{,}5$\,\%). Sur des données non attaquées, l'heuristique coupe en deux un jeu qu'un seul MR-Sort explique,
        et garde un modèle appris sur une partie seulement des données. Elles sont aussi très variables d'un modèle
        A à l'autre (intervalles de $\pm 5$ à $\pm 16$ points).
  \item \textbf{D1 ne protège pas contre S2}, l'attaque qu'elle visait : à 50\,\%, elle fait moins bien que lincs seul
        ($62{,}6$ contre $81{,}8$\,\%). Quand les deux parts du jeu sont de même taille, garder le plus grand segment
        revient à choisir presque au hasard entre A et le modèle de C.
  \item \textbf{Chaque défense spécialisée ne gagne que dans quelques cas}, et souvent sans écart significatif : D3
        contre S4 et S5 à 50\,\%, D1 contre S6 à $1/3$, D2 contre S3 et S6 à 50\,\%. Aucune défense ne domine les
        autres partout, d'où l'idée de la phase~2 : détecter l'attaque, puis choisir la défense.
\end{itemize}
Avec 20 modèles par case, les écarts de moins de 5 à 10 points sont rarement significatifs. La colonne D0 refait
l'expérience de la partie~1 avec d'autres modèles A : elle la retrouve à 4 points près, sauf contre S2 à 50\,\%
($81{,}8$ contre $72{,}2$\,\%), un écart qui rappelle que ces moyennes varient encore d'un tirage à l'autre.

\subsection{Phase 2 : détecter l'attaque, puis choisir la défense}

\paragraph{Le détecteur.} B calcule sur le jeu reçu quelques indicateurs : $n$ et $p$ ; la part des objets
impliqués dans une violation de monotonie, et la plus grande part des violations portée par un même objet ; la
part de la catégorie la plus fréquente et celle de la catégorie extrême la plus fréquente ; la concentration des
points au centre de l'espace (part des points dans $[0{,}25\,;0{,}75]^n$, rapportée à celle de points
uniformes) ; la meilleure précision d'un classement sur un seul critère ; et l'ajustement d'un MR-Sort appris par
lincs sur tout le jeu. Un arbre de décision (profondeur 5, au moins 5 jeux par feuille) prédit à partir de ces
indicateurs l'une des 13 situations de la phase~1 : pas d'attaque, ou l'une des six stratégies à $1/3$ ou à
50\,\%. L'arbre est appris \textbf{une seule fois}, sur des jeux que B simule lui-même en tirant des modèles A au
hasard et en jouant chaque attaque (30 jeux par stratégie et par proportion, 60 sans attaque). Il prédit la
bonne situation sur 53\,\% de ces jeux, et sur 42\,\% en validation croisée à 5 plis.

\paragraph{La redirection.} B applique ensuite la défense qui avait le meilleur accord moyen dans la phase~1 pour
la situation prédite (les cases en gras du tableau~\ref{tab:defenses}) ; s'il prédit qu'il n'y a pas d'attaque, il
applique D2, à égalité avec D0 sur cette ligne. La défense hybride est évaluée sur 60 nouveaux modèles A par case,
tirés indépendamment de ceux qui ont servi à apprendre l'arbre et à choisir la redirection.

\begin{table}[h]
\centering
\small
\begin{tabularx}{\textwidth}{>{\raggedright\arraybackslash}p{3.3cm}cCCCc}
\toprule
Attaque & C & Hybride (60) & lincs seul, partie 1 (60) & D2 (20) & Situation reconnue \\
\midrule
Sans attaque         & 0       & $98{,}2 \pm 0{,}3$ & $97{,}5 \pm 1{,}0$ & $98{,}6 \pm 0{,}4$ & 56/60 \\
\midrule
S1 -- Bruit          & $1/3$   & $96{,}0 \pm 0{,}7$ & $95{,}5 \pm 1{,}0$ & $96{,}7 \pm 1{,}1$ & 31/60 \\
                     & 50\,\%  & $91{,}7 \pm 2{,}2$ & $94{,}5 \pm 1{,}3$ & $93{,}3 \pm 2{,}0$ & 0/60 \\
S2 -- Autre MR-Sort  & $1/3$   & $94{,}0 \pm 0{,}9$ & $94{,}0 \pm 1{,}3$ & $94{,}7 \pm 1{,}6$ & 5/60 \\
                     & 50\,\%  & $72{,}8 \pm 4{,}5$ & $72{,}2 \pm 4{,}5$ & $74{,}0 \pm 4{,}6$ & 17/60 \\
S3 -- Extrême        & $1/3$   & $90{,}8 \pm 1{,}8$ & $92{,}3 \pm 1{,}7$ & $90{,}9 \pm 3{,}3$ & 0/60 \\
                     & 50\,\%  & $64{,}8 \pm 4{,}9$ & $65{,}4 \pm 6{,}0$ & $70{,}0 \pm 7{,}8$ & 29/60 \\
S4 -- Dictature      & $1/3$   & $92{,}3 \pm 1{,}4$ & $92{,}5 \pm 1{,}4$ & $91{,}5 \pm 3{,}1$ & 25/60 \\
                     & 50\,\%  & $70{,}7 \pm 4{,}6$ & $68{,}1 \pm 4{,}0$ & $65{,}1 \pm 5{,}5$ & 0/60 \\
S5 -- Anti-monotone  & $1/3$   & $93{,}4 \pm 1{,}1$ & $94{,}5 \pm 1{,}2$ & $94{,}2 \pm 2{,}2$ & 33/60 \\
                     & 50\,\%  & $83{,}9 \pm 4{,}1$ & $85{,}9 \pm 3{,}8$ & $76{,}3 \pm 8{,}2$ & 30/60 \\
S6 -- Extrême centre & $1/3$   & $76{,}8 \pm 4{,}8$ & $79{,}7 \pm 3{,}7$ & $72{,}4 \pm 8{,}8$ & 45/60 \\
                     & 50\,\%  & $69{,}0 \pm 4{,}2$ & $70{,}6 \pm 4{,}2$ & $77{,}0 \pm 7{,}3$ & 49/60 \\
\midrule
Pire cas sous attaque &        & $64{,}8$ & $65{,}4$ & $65{,}1$ & \\
\bottomrule
\end{tabularx}
\caption{Accord moyen (\%) entre B et A pour la défense hybride, comparée à lincs seul (partie~1) et au nettoyage
de la monotonie (phase~1), avec l'intervalle de confiance à 95\,\% (entre parenthèses, le nombre de modèles A par
case). Dernière colonne : nombre de jeux pour lesquels l'arbre a reconnu la bonne attaque et la bonne
proportion.}
\label{tab:hybrid}
\end{table}

Le tableau~\ref{tab:hybrid} conduit aux constats suivants :
\begin{itemize}
  \item \textbf{lincs seul et D2 sont déjà de très bonnes défenses.} Sur la plupart des cases, ils font mieux que les
        défenses à deux modèles, et sur les autres ils en sont très proches (tableau~\ref{tab:defenses}) : il reste
        peu de marge à une défense plus élaborée.
  \item \textbf{La défense hybride est un bon compromis entre les deux.} Elle est aussi bonne que lincs seul et D2 sur
        la plupart des cases, parfois un peu meilleure (contre S4 à 50\,\% : $70{,}7$\,\% contre $68{,}1$ et
        $65{,}1$\,\%), parfois un peu moins bonne (contre S1 à 50\,\% : $91{,}7$ contre $94{,}5$\,\%), sans écart
        significatif. Son pire cas ($64{,}8$\,\%, contre S3 à 50\,\%) est celui de lincs seul et de D2. Elle ne coûte
        presque rien sans attaque ($98{,}2$\,\%) et, surtout, n'a jamais l'effondrement des défenses à deux modèles
        appliquées à l'aveugle.
  \item \textbf{Le détecteur reconnaît bien les attaques qui laissent une trace.} L'absence d'attaque est reconnue
        dans 56 cas sur 60, et S6 dans tous les cas (la concentration au centre est visible), avec la bonne proportion
        dans 75 à 80\,\% des cas. Le premier critère de l'arbre est l'ajustement de lincs sur le jeu reçu : un jeu qu'un
        seul MR-Sort explique parfaitement est déclaré sain.
  \item \textbf{Les attaques mal reconnues sont surtout celles qui font peu de dégâts} : S3 à $1/3$ (0 sur 60), S2 à
        $1/3$ (5 sur 60) ou le bruit, que lincs écarte de toute façon. Une erreur de détection coûte alors peu, car
        la défense choisie à tort est le plus souvent D0 ou D2.
\end{itemize}
La limite de cette approche est celle de la phase~1 : la redirection ne peut pas faire mieux que la meilleure
défense de chaque situation, et contre S3 et S6 à 50\,\% aucune défense ne remonte nettement l'accord. Ces deux
attaques restent les points faibles de B, et c'est par là que C peut répondre.

\subsection{Phase 3 : la réponse de C}

C attaque toujours à l'aveugle (il ne connaît que $n$ et $p$), mais sait que B peut se défendre. Il cherche une
stratégie qui batte la défense hybride tout en restant efficace contre les autres défenses. Nous partons de S3,
la meilleure attaque à 50\,\% : tous les points de C sont dans une catégorie extrême $e$, la pire ou la meilleure,
tirée au hasard. Au lieu de les répartir uniformément, C les place là où cette étiquette contredit le plus A : si
$e$ est la pire catégorie, dans une boîte du haut de l'espace, là où tout MR-Sort classe haut ; si $e$ est la
meilleure, dans la boîte symétrique du bas. Nous comparons cinq placements, pour $e$ la pire catégorie (le cas
$e$ meilleure est symétrique) : $[0\,;1]^n$ (S3), $[0{,}25\,;1]^n$ (décalée), $[0{,}5\,;1]^n$ (haute),
$[0{,}4\,;0{,}9]^n$ (S6 décalée) et une moitié des points uniformes, l'autre au centre (mélange de S3 et S6).

\begin{table}[h]
\centering
\small
\begin{tabularx}{\textwidth}{>{\raggedright\arraybackslash}p{4.2cm}cCCCc}
\toprule
Placement des points & C & lincs seul & D2 & Hybride & Pire cas pour C \\
\midrule
S3 : $[0\,;1]^n$                 & $1/3$  & $91{,}1 \pm 3{,}0$  & $88{,}9 \pm 3{,}2$  & $91{,}7 \pm 3{,}2$  & $91{,}7$ \\
                                 & 50\,\% & $59{,}4 \pm 10{,}7$ & $57{,}2 \pm 7{,}5$  & $60{,}1 \pm 8{,}7$  & $60{,}1$ \\
Décalée : $[0{,}25\,;1]^n$       & $1/3$  & $90{,}5 \pm 5{,}0$  & $83{,}4 \pm 6{,}8$  & $80{,}9 \pm 6{,}8$  & $90{,}5$ \\
                                 & 50\,\% & $43{,}4 \pm 13{,}5$ & $47{,}2 \pm 11{,}8$ & $42{,}6 \pm 14{,}0$ & $47{,}2$ \\
\textbf{Haute : $[0{,}5\,;1]^n$} & $1/3$  & $82{,}4 \pm 9{,}9$  & $80{,}7 \pm 7{,}8$  & $70{,}1 \pm 10{,}2$ & $\mathbf{82{,}4}$ \\
                                 & 50\,\% & $30{,}8 \pm 13{,}5$ & $47{,}0 \pm 15{,}2$ & $17{,}6 \pm 10{,}3$ & $\mathbf{47{,}0}$ \\
S6 décalée : $[0{,}4\,;0{,}9]^n$ & $1/3$  & $83{,}5 \pm 7{,}1$  & $63{,}9 \pm 10{,}0$ & $76{,}9 \pm 10{,}8$ & $83{,}5$ \\
                                 & 50\,\% & $52{,}4 \pm 10{,}4$ & $58{,}0 \pm 11{,}2$ & $42{,}9 \pm 9{,}5$  & $58{,}0$ \\
Mélange S3 et S6                 & $1/3$  & $84{,}7 \pm 5{,}2$  & $76{,}2 \pm 5{,}4$  & $78{,}9 \pm 8{,}8$  & $84{,}7$ \\
                                 & 50\,\% & $62{,}6 \pm 7{,}6$  & $69{,}2 \pm 9{,}4$  & $59{,}1 \pm 7{,}8$  & $69{,}2$ \\
\bottomrule
\end{tabularx}
\caption{Accord moyen (\%) entre B et A pour les variantes de S3, avec l'intervalle de confiance à 95\,\% (20
modèles A tirés indépendamment pour chaque case). Dernière colonne : l'accord contre la défense qui résiste le
mieux.}
\label{tab:counter}
\end{table}

La variante \textbf{haute} est la meilleure aux deux proportions (tableau~\ref{tab:counter}) : elle fait mieux que
S3 contre les trois défenses, et son pire cas passe de $91{,}7$ à $82{,}4$\,\% à $1/3$, et de $60{,}1$ à $47{,}0$\,\% à
50\,\%. Elle exploite la monotonie de MR-Sort : un MR-Sort ne peut pas classer le haut de l'espace dans la pire
catégorie sans y classer tout l'espace. Quand C est assez nombreux, lincs préfère donc tout mettre dans la pire
catégorie, et le modèle appris s'effondre. Elle échappe aussi au détecteur : ses points ne sont pas concentrés au
centre, et l'arbre, qui ne l'a jamais vue, la prend pour S3 et la renvoie vers D2 dans les 20 cas sur 20 à
50\,\%. L'écart entre l'hybride ($17{,}6$\,\%) et D2 ($47{,}0$\,\%) à 50\,\% vient donc du tirage des modèles A, et
non de la défense : l'attaque est en tout ou rien, et selon le modèle A, lincs s'effondre complètement ou
résiste.

Contre les défenses à deux modèles, la variante haute bat D1 ($80{,}6 \pm 6{,}8$\,\% à $1/3$, $44{,}8 \pm 15{,}6$\,\% à
50\,\%) et D3 à $1/3$ ($81{,}5 \pm 8{,}8$\,\%), mais \textbf{D3 lui résiste à 50\,\%} ($79{,}5 \pm 9{,}6$\,\%, contre
$66{,}9$\,\% face à S3). Les points de C étant tous dans la même catégorie, l'heuristique à deux modèles les isole dans
un segment qu'elle explique par un modèle dégénéré, où un seul critère atteint le seuil de majorité : D3 le prend
pour un dictateur et le rejette. C'est un effet de bord de la définition de D3, pas une détection de l'attaque.
Nous avons essayé de rendre les données de C non dictatoriales (une partie des points dans le coin opposé,
étiquetée à l'autre extrême), sans succès : l'heuristique apprend encore un modèle dictatorial pour C dans la
moitié des cas.

Nous retenons la variante haute comme réponse de C. Il n'existe pas d'attaque meilleure que les autres contre
toutes les défenses, ni de défense meilleure contre toutes les attaques : améliorer l'une contre un type de
stratégie a tendance à la dégrader contre d'autres. Les variantes de S3 restent cependant globalement les
meilleures attaques, car elles sont très difficiles à défendre ; leur seule parade, D3, coûte environ 15 points
à B lorsqu'il n'y a pas d'attaque.

% =====================================================================
\section{Quatrième scénario : l'attaquant a accès aux données de A, B se défend}

Nous reprenons le cadre de la partie~2 (C reçoit les données de A et en infère $\hat{A}$) face à la défense
hybride de la partie~3, en trois phases : les attaques informées de la partie~2 contre l'hybride ; une version
informée de S3 ; enfin l'attaque la plus prometteuse contre toutes les défenses, qui constitue notre réponse
finale.

\subsection{Phase 1 : les attaques informées contre la défense hybride}

\begin{table}[h]
\centering
\begin{tabularx}{\textwidth}{>{\raggedright\arraybackslash}p{5.0cm}cCC}
\toprule
Stratégie & C & Hybride & lincs seul (partie~2) \\
\midrule
S7 -- Modèle $\hat{A}$ décalé       & $1/3$  & $86{,}2 \pm 1{,}8$ & $89{,}8 \pm 1{,}9$ \\
                                    & 50\,\% & $56{,}1 \pm 4{,}3$ & $55{,}0 \pm 2{,}7$ \\
S8 -- Copie à étiquette décalée     & $1/3$  & $91{,}7 \pm 2{,}9$ & $94{,}2 \pm 1{,}0$ \\
                                    & 50\,\% & $52{,}5 \pm 5{,}0$ & $51{,}5 \pm 4{,}1$ \\
S9 -- Viser les frontières          & $1/3$  & $86{,}2 \pm 1{,}7$ & $87{,}5 \pm 1{,}6$ \\
                                    & 50\,\% & $77{,}5 \pm 3{,}3$ & $79{,}3 \pm 2{,}3$ \\
S10 -- Optimisation (10 modèles)    & $1/3$  & $\mathbf{62{,}0 \pm 6{,}7}$ & $62{,}1 \pm 9{,}1$ \\
                                    & 50\,\% & $\mathbf{33{,}2 \pm 7{,}3}$ & $23{,}1 \pm 13{,}8$ \\
\bottomrule
\end{tabularx}
\caption{Accord moyen (\%) entre B et A pour les attaques informées, contre la défense hybride et contre lincs
seul, avec l'intervalle de confiance à 95\,\% (60 modèles A tirés indépendamment par case, 10 pour S10).}
\label{tab:informed-hybrid}
\end{table}

La défense hybride ne protège pas mieux que lincs seul contre les attaques informées
(tableau~\ref{tab:informed-hybrid}) : tous les écarts restent dans les intervalles de confiance, et la plupart sont
même en faveur de C à $1/3$. Le détecteur n'a jamais vu ces attaques ; il les prend le plus souvent pour S5 ou S3, et
les renvoie vers D2, D3 ou D0, qui ne les contrent pas. S10, qui simule pourtant une procédure inv-MR-Sort générique et non la défense de B, reste la
meilleure attaque informée. À 50\,\%, elle ne fait pas mieux que la variante haute, qui attaque à l'aveugle
($17{,}6$\,\%).

\subsection{Phase 2 : S3 informée}

C place ses points comme dans S3, tous dans une catégorie extrême, uniformément dans une boîte
$[c-r,\,c+r]$, mais choisit la boîte et la catégorie grâce aux données de A. Un algorithme génétique, réglé comme
pour S10 (10 individus, 7 générations), optimise le centre $c$, la demi-largeur $r$ et la catégorie (pire ou
meilleure) ; chaque individu est évalué en simulant une procédure inv-MR-Sort générique (lincs) et en mesurant
l'accord entre le modèle obtenu et $\hat{A}$. C ne simule pas la défense hybride : dans la réalité, l'attaquant ne
sait pas comment le défenseur se défend, et le supposer serait une hypothèse forte ; il peut en revanche toujours
simuler une procédure inv-MR-Sort générique. La population de départ contient S3, S6 et la variante haute,
chacune dans les deux sens. L'attaque est évaluée contre la défense hybride sur 10 modèles A par proportion.

\begin{table}[h]
\centering
\begin{tabularx}{0.8\textwidth}{>{\raggedright\arraybackslash}p{5.5cm}CC}
\toprule
Attaque contre l'hybride & $1/3$ & 50\,\% \\
\midrule
S3 informée (10 modèles)                   & $69{,}7 \pm 14{,}6$ & $\mathbf{19{,}5 \pm 8{,}9}$ \\
S10 (10 modèles)                           & $\mathbf{62{,}0 \pm 6{,}7}$ & $33{,}2 \pm 7{,}3$ \\
Variante haute, à l'aveugle (20 modèles)   & $70{,}1 \pm 10{,}2$ & $17{,}6 \pm 10{,}3$ \\
\bottomrule
\end{tabularx}
\caption{Accord moyen (\%) entre B, protégé par la défense hybride, et A, pour les trois meilleures attaques.}
\label{tab:informed-s3}
\end{table}

Les trois attaques sont proches, et leurs écarts restent dans les intervalles de confiance
(tableau~\ref{tab:informed-s3}). L'algorithme génétique \textbf{retrouve de lui-même l'idée de la variante haute} :
dans 15 cas sur 20, il choisit une boîte dans le bas de l'espace (centre moyen entre $0{,}07$ et $0{,}43$), étiquetée
dans la meilleure catégorie, et l'élargit à 50\,\% (demi-largeur jusqu'à $0{,}47$). À $1/3$, l'attaque réussit
(accord de 38 à 67\,\%) quand l'hybride renvoie le jeu vers D2, et échoue (93 à 99\,\%) dans les 3 cas où il le
renvoie vers D1 ou D3, d'où la largeur de l'intervalle.

\subsection{Phase 3 : la réponse finale de C}

Nous retenons la S3 informée, la meilleure attaque à 50\,\% contre l'hybride, qui généralise la variante haute en
s'adaptant à chaque modèle A, et l'évaluons contre toutes les défenses (10 modèles A par case).

\begin{table}[h]
\centering
\begin{tabularx}{0.8\textwidth}{>{\raggedright\arraybackslash}p{5.5cm}CC}
\toprule
Défense de B & $1/3$ & 50\,\% \\
\midrule
D0 -- lincs seul                & $54{,}9 \pm 11{,}0$ & $\mathbf{6{,}1 \pm 3{,}8}$ \\
D2 -- Nettoyage de la monotonie & $\mathbf{53{,}3 \pm 6{,}4}$ & $22{,}7 \pm 11{,}7$ \\
D1 -- Plus grand segment        & $69{,}5 \pm 16{,}7$ & $9{,}0 \pm 7{,}9$ \\
Hybride (phase~2)               & $69{,}7 \pm 14{,}6$ & $19{,}5 \pm 8{,}9$ \\
D3 -- Rejet du dictateur        & $75{,}3 \pm 15{,}5$ & $89{,}4 \pm 16{,}0$ \\
\bottomrule
\end{tabularx}
\caption{Accord moyen (\%) entre B et A face à la S3 informée, pour chaque défense, avec l'intervalle de confiance
à 95\,\% (10 modèles A tirés indépendamment par case).}
\label{tab:final}
\end{table}

Contre lincs seul, D2, D1 et l'hybride, la S3 informée est l'attaque la plus forte de ce travail
(tableau~\ref{tab:final}) : à 50\,\%, l'accord tombe entre 6 et 23\,\%, et à $1/3$ entre 53 et 70\,\%. Sa seule parade
est D3, qui la neutralise presque à 50\,\% (plus de 93\,\% d'accord dans 9 cas sur 10), pour la même raison qu'avec
la variante haute : à parts égales, l'heuristique à deux modèles sépare proprement les données de C, toutes dans la
même catégorie, de celles de A, les explique par un modèle dégénéré que D3 rejette, et B garde le modèle appris sur
les seules données de A. Mais cette parade est un mauvais choix pour B : D3 coûte environ 15 points lorsqu'il n'y
a pas d'attaque ($83{,}5$\,\%), reste moins bonne que lincs seul contre la plupart des attaques à l'aveugle
(tableau~\ref{tab:defenses}), et B ne peut pas savoir à l'avance qu'il fait face à cette attaque : le détecteur ne
la lui a presque jamais désignée. \textbf{La S3 informée est notre réponse finale.}

% =====================================================================
\section*{Conclusion}

\paragraph{Attaquer.} Un attaquant minoritaire n'a presque aucun effet sur un apprentissage inv-MR-Sort, tant qu'il
répartit ses données : lincs maximise le nombre d'objets bien classés et écarte des contradictions dispersées. Ce
qui compte est la \emph{proportion locale} : en concentrant ses points dans une zone bien choisie, C y devient
majoritaire alors qu'il reste minoritaire globalement. La structure de MR-Sort fournit cette zone : sa monotonie
empêche de classer le haut de l'espace dans la pire catégorie sans y classer tout l'espace, si bien que des points
étiquetés à l'extrême opposé de ce que tout MR-Sort leur attribue forcent l'apprentissage à choisir entre A et C. Les
attaques les plus efficaces envoient ainsi des données cohérentes, représentables par un MR-Sort, mais
contraires au modèle de A. Connaître les données de A et pouvoir simuler une procédure inv-MR-Sort générique permet d'adapter ce
placement à chaque modèle A : l'optimisation retrouve d'elle-même la variante haute, et fait tomber l'accord de B
avec A à quelques pourcents lorsque C fournit la moitié des données.

\paragraph{Se défendre.} lincs seul et le nettoyage de la monotonie sont déjà de très bonnes défenses contre les
attaques simples : ils écartent le bruit et les données incohérentes sans rien coûter lorsqu'il n'y a pas
d'attaque. Les défenses fondées sur deux modèles coûtent cher en l'absence d'attaque et sont très variables. La
défense hybride, qui détecte l'attaque puis choisit la défense, est un bon compromis : aussi bonne que les meilleures
défenses sur la plupart des situations, sans leur effondrement. Mais elle ne fait pas mieux que la meilleure
défense de chaque situation, et ne reconnaît que les attaques qu'elle a vues pendant son apprentissage.

\paragraph{Un jeu de compromis.} Aucune attaque n'est meilleure que les autres contre toutes les défenses, et aucune
défense n'est meilleure contre toutes les attaques : améliorer l'une contre un type de stratégie la dégrade contre
d'autres. La variante haute et la S3 informée battent toutes les défenses sauf D3, qui les contre par un effet de
bord, mais que B a toutes les raisons de ne pas choisir. Les variantes de S3 sont globalement les meilleures
attaques, car très difficiles à défendre : face à des données cohérentes mais contraires à A, un apprenant qui ne
sait pas quelle part de ses données est fiable n'a pas de moyen sûr de les écarter.

\paragraph{Limites.} Nos conclusions reposent sur 10 à 60 modèles A par situation ; avec 10 ou 20 modèles, les
intervalles de confiance restent larges, et beaucoup d'écarts de moins de 10 points ne sont pas significatifs. Nous
n'avons étudié que deux proportions pour les défenses ($1/3$ et 50\,\%), une seule procédure d'apprentissage (lincs)
et des données de A non bruitées. Enfin, le tableau~\ref{tab:algo} de la partie~1 a été obtenu avant la
correction de la section~0 et n'a pas été recalculé.

\end{document}
